\cotefile{46}
\scanpage{49}
% archive sheet number: 48
% Faint reversed typescript is show-through from the other side.
\[
\underline{N}\simeq\underline{L}\otimes\unsure{\mathscr{O}|\mathcal{O}}_{\widetilde{P}}(1)
\]
\[
\underline{N}\otimes_{\mathscr{O}_{S}}\underline{L}^{-1}=\unsure{M|\mathscr{M}}
\quad\text{\unsure{module} \unsure{inversible} sur }X\times_{S}\widetilde{P}.
\]

Loc. au dessus de $S$ il \unsure{provient} d'un module inversible sur $\widetilde{P}$, d'ailleurs \uline{de façon canonique}. \unsure{grâce au} fait que $f_{*}(\mathscr{O}_{\deleted{\unsure{P}}X})=\mathscr{O}_{S}$ [\unsure{pour le} \unsure{préciser|principe} que ce soit \unsure{vrai} \unsure{univ.},] $f$ étant \unsure{supposé} \unsure{qu. cpt} \unsure{qu. sép}].
\scanpage{50}
\mspage{- 59 -}
% archive sheet number: 49
n\handwritten{°357}

p-adique $\handwritten{|}x\handwritten{|}_{p}=p^{-\unsure{v|\nu}_{p}(x)}$ (\uline{Top.Gén.} Chap.IX,\handwritten{§3}, n\handwritten{°2}) ; soit $\mathbb{Z}_{p}$ l'anneau des entiers p-adiques, adhérence de $\mathbb{Z}$ dans $\mathbb{Q}_{p}$ ; $\mathbb{Z}_{p}$ est un sous-groupe du groupe additif de $\mathbb{Q}_{p}$, et c'est aussi la boule ouverte $\handwritten{|}x\handwritten{|}_{p}\handwritten{<}3/2$ ; c'est donc un sous-groupe ouvert et fermé de $\mathbb{Q}_{p}$, de sorte que $\mathbb{Q}_{p}\handwritten{/}\mathbb{Z}_{p}$ est discret. De plus, $\mathbb{Z}_{p}$ est compact, par exemple parce qu'on peut l'identifier d'une manière évidente avec la limite projective des groupes $\mathbb{Z}\handwritten{/}p^{n}\mathbb{Z}$. Il est clair qu'il y a, sur le groupe additif $\mathbb{Q}_{p}$, une mesure de Haar $\handwritten{\alpha}$ et une seule telle que $\handwritten{\alpha}(\mathbb{Z}_{p})=1$ ; celle-ci est dite la \uline{mesure de Haar normalisée} sur $\mathbb{Q}_{p}$. Comme $\mathbb{Z}_{p}$ se compose, pour tout $n\handwritten{\geq}0$, de $p^{n}$ classes suivant le sous-groupe (ouvert et compact) $p^{-n}\mathbb{Z}_{p}$, on a $\handwritten{\alpha}(p^{n}\mathbb{Z}_{p})=p^{-n}$ ; on a de même $\handwritten{\alpha}(p^{-n}\mathbb{Z}_{p})=p^{n}$, puisque $p^{-n}\mathbb{Z}_{p}$, pour $n\handwritten{\geq}0$, se compose de $p^{n}$ classes suivant $\mathbb{Z}_{p}$. Soient $A_{1},\ldots,A_{m}$ des classes suivant $p^{n}\mathbb{Z}_{p}$ ; soit $f=\handwritten{\sum}_{i=1}^{m}a_{i}\handwritten{\varphi}_{A_{i}}$ ; on aura donc $\handwritten{\int}f\mathrm{d}\handwritten{\alpha}=p^{-n}\handwritten{\sum}_{i}a_{i}$. Soit maintenant $g\handwritten{\in}\handwritten{\unsure{\mathscr{K}|\mathcal{K}}}(\mathbb{Q}_{p})$ ; soit $K$ une réunion finie de classes suivant $\mathbb{Z}_{p}$, contenant le support de $g$. Comme $g$ est uniformément continue, il y a, pour tout $\handwritten{\varepsilon> }0$, un $n\handwritten{\geq}0$ tel que $\handwritten{|}g(x)-g(y)\handwritten{|\leq\varepsilon}$ pour $\handwritten{|}x-y\handwritten{|}_{p}\handwritten{\leq}p^{-n}$, c'est-à-dire pour $x\handwritten{\equiv}y\ \mathrm{mod.}\ p^{n}$ ; soient $A_{1},\ldots,A_{m}$ les classes suivant $p^{n}\mathbb{Z}_{p}$ contenues dans $K$ ; soient $a_{1},\ldots,a_{m}$ des représentants de ces classes ; soit $f=\handwritten{\sum}_{i=1}^{m}g(a_{i})\handwritten{\varphi}_{A_{i}}$ ; on a alors $\handwritten{|}f-g\handwritten{|\leq\varepsilon}$. En particulier, si $f$ ne prend qu'un nombre fini de valeurs distinctes $c_{\unsure{o|0}}=0$, $c_{1},\ldots,c_{r}$, et qu'on prenne $\handwritten{\varepsilon<}\inf_{i\handwritten{\ne}j}\handwritten{|}c_{i}-c_{j}\handwritten{|}$, on a $f=g$ ; plus particulièrement, si $f$ est la fonction caractéristique $\handwritten{\varphi}_{X}$ d'une ensemble $X$, et que $\handwritten{\varphi}_{X}\handwritten{\in\mathscr{K}}(\mathbb{Q}_{p})$, ce qui revient à dire que $X$ est ouvert et compact, on voit qu'il y a $n\handwritten{\geq}0$ tel que $X$ soit réunion des classes suivant $p^{n}\mathbb{Z}_{p}$ qui y sont contenues. Soit $\handwritten{\unsure{\Phi|\mathfrak{P}}}$ la famille formée de telles réunions, c'est-à-dire des parties ouvertes et compactes de $\mathbb{Q}_{p}$ ; c'est un clan au sens du Chap.IV,\handwritten{§4}, n\handwritten{°8}, et les fonctions $\handwritten{\Phi}$-étagées sont, d'après ce qu'on vient de voir, les fonctions
\scanpage{51}
% archive sheet number: 50
Progroupes et \unsure{espaces} \unsure{homogènes} \unsure{quasi-compacts|quasi-compact}
\scanpage{52}
% archive sheet number: 51
\margin{\uline{! : \unsure{cas.} affine} \deleted{au \unsure{sens} \unsure{syst. proj.} de \uuline{\unsure{groupes}}}}
\diagram{\begin{tikzcd}[ampersand replacement=\&,column sep=large,row sep=large]
\operatorname{Progr}^{!}(C'_0)
  \arrow[r,"\deleted{\unsure{\leftarrow}}","\alpha"']
\& \operatorname{Gr}(C)
  \arrow[r,phantom,"{\underset{\beta}{\overset{\text{\unsure{équiv.}}}{\simeq}}}" description]
\& \operatorname{Gr}(C')
  \arrow[r,"{\substack{\text{pl. fid}\\\deleted{\unsure{\operatorname{Gr}(C)}}}}","\unsure{\sigma|\tau}"']
\& \operatorname{Gr}\ \text{qucpt}\ (\mathrm{Sch}_{/k})
  \arrow[r,"{\substack{\deleted{\text{pl. fid}}\\\unsure{\ell}}}"]
\& {} \\
\operatorname{Progr}(C'_0)\arrow[u,hook,"i\ \text{pl. fid}"]
  \arrow[rr,"\unsure{\kappa|k}"']
\& \& \operatorname{Gr}\operatorname{Pro}(C'_0)
  \arrow[u,hook,"j\ \text{pl. fid}"']
\end{tikzcd}}

$\beta$ \unsure{équivalence} car $\deleted{\unsure{\operatorname{Gr}(C)\simeq\operatorname{Gr}(}}\ C\simeq C'$.

$\sigma$ pl. fid \unsure{trivial} sur $\operatorname{Gr}(C)\longrightarrow\operatorname{Gr}\ \text{qu. cpts}\ (\mathrm{Sch}_{/k})$.

$i,j$ pl. fid \deleted{\unsure{est}} \unsure{trivial} [N.B $C'\unsure{-|\longrightarrow}\operatorname{Pro}(C'_0)$ est pl. fid et \unsure{commute} aux $\varprojlim$ \unsure{finies}]

\[
\kappa\ \text{pl. fid}\ \Longrightarrow\ \alpha\ \text{pl. fid}\ \Longleftrightarrow\
[\sigma\alpha:\operatorname{Progr}^{!}(C'_0)\longrightarrow\operatorname{Gr}\ \text{qucpt}\ (\mathrm{Sch}_{/k})\ \deleted{\text{\unsure{est}}}\ \text{pl. fid.}]
\]
\margin{\boxed{\text{vrai}}}
% The boxed "vrai" is joined by a line to the circled assertion kappa pl. fid.
\diagram{\begin{tikzcd}[ampersand replacement=\&,column sep=large,row sep=large]
\sigma\alpha\ \text{une \unsure{équivalence}}
  \arrow[r,Leftrightarrow]
\& \sigma\ \text{et}\ \alpha\ \text{des \unsure{équiv.}}
  \arrow[d,Rightarrow]\arrow[r,Rightarrow]
\& \begin{gathered}\deleted{\text{\unsure{tout}}}\ \text{schéma en}\\
\text{\unsure{gr.} qu. cpt sur}\\ k\ \text{admet un}\\
\text{\unsure{noyau}}\ \text{\uline{affine}}\end{gathered} \\
\kappa\ \text{une \unsure{équivalence}}\arrow[r,Rightarrow]
\& \alpha\ \text{une \unsure{équivalence}}
\& \begin{gathered}\text{donc un schéma en}\\
\text{\unsure{gr. quot} de t. f. sur }k\end{gathered}
  \arrow[ul,Leftarrow,bend left=25,"\text{\unsure{un peu de travail}}"']
\end{tikzcd}}
\margin{$\unsure{\kappa|K}$}
\margin{\unsure{douteux}}
% The circled "douteux" is linked to "kappa une équivalence".
\margin{\unsure{pourvu qu'on} \deleted{\unsure{sache}} \unsure{qu'un|puisse} \unsure{un} \unsure{pro. gr.} tel que le \unsure{proschéma} \unsure{associé} \unsure{soit} \unsure{un.} affine, \unsure{soit} \unsure{lui-même} \unsure{un} \unsure{sch.} affine.}
% This last note is connected to the implication kappa equivalence => alpha equivalence.
\scanpage{53}
% archive sheet number: 52
Tout objet de $\operatorname{Progr}(C'_0)$ \unsure{provient} d'un système \uline{projectif strict}. \unsure{Considérons} donc deux $(G_i)_i$, $(H_j)_j$ et \unsure{montrons} que
\[
\unsure{\varphi|\phi}:\ \varprojlim_j\varinjlim_i\operatorname{Hom}_{\mathrm{gr}}(G_i,H_j)
\unsure{-|\longrightarrow}
\left(\varprojlim_j\varinjlim_i\operatorname{Hom}(G_i,H_j)\right)^{\mathrm{mult.}}
\]
est \uline{bijectif}. En effet, comme pour $i,j$ fixés,
\[
\operatorname{Hom}_{\mathrm{gr}}(G_i,H_j)-\operatorname{Hom}(G_i,H_j)
\]
est \unsure{injectif}, on voit que $\varphi$ est \uline{injectif}. \deleted{\unsure{Pour le}}

Montrons que c'est \uline{surjectif}, \deleted{\unsure{Prouvons|Revenons}} \inserted{donc} que \unsure{si|à} $(u_j)\in\varprojlim_j\varinjlim_i\operatorname{Hom}(G_i,H_j)$ est multiplicatif, \unsure{tous les}
\margin{$\unsure{G_i}$}
\[
u_j:(G_i)_i-H_j
\]
\deleted{\unsure{proviennent}} d'un $u_j^{(i)}\ G_i-H_j$ \uuline{mult.} \deleted{\unsure{Car}}
[si c'est vrai pour un, c'est vrai pour tous \unsure{un}].
\diagram{\begin{tikzcd}[ampersand replacement=\&,column sep=large,row sep=large]
G_j\times G_j\arrow[dr]\&\deleted{\unsure{G_j\times G_j}}\&{}\\
\&G_i\times G_i\arrow[r]\arrow[d]\&H\times H\arrow[d]\\
\&\deleted{\unsure{C}}G_i\deleted{\unsure{G}}\arrow[r,no head]\&H
\end{tikzcd}}
\unsure{diagramme} devient comm. pour $\unsure{i'\longrightarrow|i'}$ \unsure{car}
\[
G_j\times G_j-G_i\times G_i\qquad\text{pour }i'\ \text{\unsure{grand},}
\]
\deleted{$G_j\longrightarrow G_i$} \uline{est un \uuline{épi} de \unsure{schémas}}, donc l'$\unsure{\Delta}$ est \unsure{comm.} \unsure{avant}.
% The product morphism is bracketed and linked down to "est un épi de schémas".
% slanted
\margin{N.B. Ce qui compte, c'est que les \unsure{morphismes} de \unsure{transition} sont des \unsure{épi} \uuline{universels}.}
\scanpage{54}
% archive sheet number: 53
\uline{Lemme} Tout \unsure{pro-groupe} \deleted{\unsure{ess. affine}} \inserted{$(G_i)_{i\in I}$} tel que le \unsure{proschéma} \unsure{associé} \deleted{\unsure{est}}soit \deleted{\unsure{un.}} affine, est ess. affine.

\deleted{\unsure{On}} On peut \unsure{supposer} $(G_i)_{i\in I}$ système projectif strict. Par hyp. il existe un syst. projectif de schémas de \unsure{prés.} finie \inserted{ess. affines} $(X_j)_{j\in J}$ qui lui est \unsure{isom.} comme \unsure{proschéma}. On peut supposer les $X_j\unsure{-|\longrightarrow}X_0$ affines.

On a
\diagram{\begin{tikzcd}[ampersand replacement=\&,column sep=large,row sep=large]
(G_i)\arrow[r,very thick]\arrow[d,no head]\&G_{i_0}\arrow[d]\\
(X_j)\arrow[r,no head]\&X_0
\end{tikzcd}}
% slanted
\margin{\uline{\unsure{c'est vrai|c'est mieux}}}
% A long curved line links this margin note to the bottom "fid plat" annotation.
Je \uline{dis que les} $G_{i_1}\longrightarrow G_{i_0}$ \uline{sont affines} \inserted{$(i_1>i_0)$}.

En effet, $\deleted{(G_i)\longrightarrow}$ $(X_j)\simeq(G_i)\longrightarrow G_{i_1}$ se \unsure{factorise} par
\[
(X_j)-X_{j_1}-G_{\unsure{j_1|i_1}},
\]
\unsure{rendant}
\diagram{\begin{tikzcd}[ampersand replacement=\&,column sep=large,row sep=large]
(G_i)\arrow[r]\arrow[d,leftrightarrow]\&G_{i_1}\arrow[r]\&G_{i_0}\arrow[d]\\
X=(X_j)\arrow[r]\&X_{j_1}\arrow[u]\arrow[r,"Q" description]\&X_0
\end{tikzcd}}
\unsure{commutatif} le carré $Q$, et
\diagram{\begin{tikzcd}[ampersand replacement=\&,column sep=large,row sep=large]
G_{i_2}\arrow[r,very thick]\arrow[rrr,bend right=45,"\text{fid plat}"']
\&X_{j_1}\arrow[r,"\text{[\unsure{affine}]}"']\arrow[rr,bend right=25,"\text{affine}"']
\&G_{i_1}\arrow[r]\&G_{i_0}
\end{tikzcd}}
\margin{\unsure{image} \unsure{schématique} est \unsure{iso} \unsure{si} $\unsure{i_2>i_1|i_2\geq i_1}$ \unsure{grand}}
\[
/
\]
\diagram{\begin{tikzcd}[ampersand replacement=\&,column sep=large]
X=\varprojlim G_i\arrow[r,very thick,"\text{fid plat}"]\arrow[rr,bend right=25,"\text{affine}"']
\&G_{i_1}\arrow[r]\&G_{i_0}
\end{tikzcd}}
\scanpage{55}
% archive sheet number: 54
% margin note written upside down
\margin{\deleted{``\unsure{Revêtements|Revêtement}'' étales de \unsure{groupes} \unsure{infinis}}}
\scanpage{56}
% archive sheet number: 55
Schémas de Brauer-\unsure{Sévéri|Séveri}
\scanpage{57}
% archive sheet number: 56
La \uline{théorie} de \uline{Brauer}-\uline{\unsure{Sévéri|Séveri}}-\uline{\unsure{Châtelet|Chatelet}} \uline{globale}.

1. \unsure{Définition} d'un schéma de Brauer-Sévéri sur $Y$ : \deleted{\unsure{isomorphe}} $\unsure{S|\mathscr{S}}$-isomorphe à $\mathbb{P}^{r}_{Y}$. \inserted{i.e. c'est un schéma de Châtelet} [\unsure{Il est donc} propre, plat sur $Y$, ses fibres sont des Brauer-Sévéri]. Possibilité de la $S$-descente, propreté. Local-isotrivialité.

2. Fibré projectif (\uline{associé} à un faisceau loc. libre) $=$ schéma de Châtelet sur $Y$, \inserted{c'est un Br.-Sé. loc. \unsure{tr.}} muni d'un $\underline{L}$ de rang $1$ qui sur chaque fibre induit \ldots\ (base noeth.); \uline{Corollaire 1} Si $\unsure{Y|S}=\operatorname{Spec}(\unsure{O|0})$, $O$ local complet, alors Châtelet $=$ Brauer-Sévéri. \uline{Corollaire 2} Si $Y$ est non singulier, alors Br. Sévéri localement trivial $=$ \deleted{fibré} projectif $=$ Châtelet qui \deleted{\unsure{un faisceau} \unsure{e.g.} \unsure{en un} \unsure{pt} (de $Y$ \unsure{supposé connexe})} \inserted{est loc. trivial en un pt de la base}
\margin{\unsure{Construction} des diviseurs ``\unsure{linéaires}''}
\inserted{\uline{Cor. 3} S'il existe un diviseur loc. principal \unsure{relatif} comme fibré associé \deleted{\unsure{déjà}} \unsure{dans} (\unsure{sur} $\unsure{\overline{X}=X|X=X'}$).}

3. Fibré dual, \unsure{interprétation} de ses sections comme diviseurs loc. principaux sur $X$ qui sont ``\unsure{spéciaux}''. \uline{Cor. 1} Si $X$ admet une section, il est un fibré projectif [$Y$ noeth.\ldots] De façon plus précise, un $X$ \inserted{de Br.-Sév.,} avec une section, c'est \unsure{la} $=$ chose que la donnée d'un \deleted{fibré} faisceau loc. libre de rang $r$, et d'un sous faisceau loc. libre de rang $r-1$, loc facteur direct\ldots

\uline{Cor. 2} \deleted{\unsure{Autre}} Critère de \unsure{locle} trivialité : existence de sections locales\ldots

4. \uline{Th} Soit $X$ propre et plat sur $Y$. \deleted{\unsure{si}} si $f^{-1}(y)$ est un Br.-Sév., alors il existe un voisinage ouvert $U$ de $y$ t.q. $X|U$ soit Br. Sév.

\uline{Cor.} Châtelet $=$ Brauer-Sévéri.

On \unsure{prouve} d'abord qu'il y a un $U$ t.q. $X|U$ soit Châtelet, puis on \deleted{\unsure{considère}} \unsure{prouve} que Châtelet $\Rightarrow$ Br. Sév., grâce à la considération des $\operatorname{Is}_{Y}(\mathbb{P}^{r}_{Y},X)$.

\{ Relations avec les faisceaux d'algèbres simples centrales\ldots\ Groupe de Brauer global. [\unsure{Consulter} le cours de Serre\ldots] \}

\uline{3 bis} Fibrés des $n$-formes associé à $\deleted{\unsure{Y}}\unsure{X}$, ses sections sur $Y$ sont les diviseurs loc. principaux sur $X$ dont la trace sur \deleted{\unsure{toute}} fibre (ou une fibre, si $Y$ est connexe) est de degré $n$.
\scanpage{58}
% archive sheet number: 57
Soit $\underline{E}$ faisceau loc. libre de rang $r$. Alors \deleted{$\unsure{\operatorname{GP}(\underline{E})}$ \unsure{opère sur}} $P(S^n(\underline{E}))$, \inserted{et $P(S^n(\underline{E})^{\vee})=P(S^n(\unsure{\check{\underline{E}}|\underline{E}^{\vee}}))$} \unsure{ne dépend que de} $P(\underline{E})$, \deleted{\unsure{de façon précise}} et $\operatorname{GP}(\underline{E})$ opère sur $P(S^n(\underline{E}))$ \deleted{\unsure{de cette manière}} \inserted{$P(\unsure{\widetilde{S}}^{n}(\underline{E})^{\vee})$} [\unsure{hom.}
\[
\operatorname{GP}(\underline{E})\longrightarrow\operatorname{GP}(S^n_{\underline{O}_{Y}}(\underline{E}))
\]
] Faisant $\underline{E}=\underline{O}_{Y}^{r+1}$, \unsure{on trouve que} $\operatorname{GP}(r)_{Y}$ opère dans $P(S^n(O_{Y}^{r+1})^{\vee})$, \unsure{d'où} \inserted{et par suite pour tout schéma de Br.-Sév. $X$} fibré principal \unsure{hom.} sous $\operatorname{GP}(r)_{Y}$, \deleted{\unsure{un fibré}\ \unsure{associé}} un fibré \deleted{\unsure{proj.}} \inserted{associé}, $\unsure{S|\mathscr{S}}$-isomorphe à $P(S^n(\underline{O}_{Y}^{r+1}))$, qui est donc un schéma de Brauer-Sévéri. Il est loc. trivial (trivial) si $X$ \deleted{l'est, \unsure{on le désigne}} \inserted{On l'appelle le} \uline{schéma des $n$-formes de} $X$, on le \unsure{dénote} par $\Phi^n(X)$. Compatibilité aux images inverses.
\margin{\unsure{n?}\ldots}
\margin{$\unsure{\operatorname{GP}(r)_{Y}}\longrightarrow\unsure{\operatorname{GP}}(S^n_{\underline{O}_{Y}}(O_{Y}^{r+1})^{\vee})\ \deleted{\unsure{\longrightarrow}}$}

Prenons d'abord $X=P(\underline{E})$. Alors
\[
\Phi^n(X)\overset{\text{can.}}{\simeq}P(\underline{S}^{n}_{\underline{O}_{X}}(\underline{E})^{\vee}).
\]
[\unsure{démonstration} : \unsure{c'est} un foncteur, compatible aux images inverses\ldots, qui définit une \unsure{représentation} et par suite une notion de fibré associé; \unsure{dans} cette \deleted{\unsure{fibré}} notion d'associé \unsure{redonne} le foncteur].

\deleted{\unsure{Sections de} $\unsure{\Phi}$} Alors une section de $\Phi^{(n)}(X)$ est donc définie par \deleted{\unsure{un}} la donnée d'un \deleted{\unsure{épi}}\inserted{homo.}\unsure{morphisme} de faisceaux de $\underline{S}^{n}(\underline{E})^{\vee}$ sur un \deleted{\unsure{loc. lib.}} faisceau inversible, \unsure{donc correspondant} à un hom ``ne s'annulant nulle part'' d'un fibré inversible dans $\underline{S}^{n}(\underline{E})$, i.e. une section de
\[
\deleted{\unsure{\mathscr{H}}}(\underline{S}^{n}(\underline{E}))^{*}/\underline{O}_{Y}^{*};
\]
\inserted{faisceau des sections de $\underline{S}^{n}(\underline{E})$ qui ne ``s'annulent pas''}
On \unsure{aura}
\scanpage{59}
% archive sheet number: 58
\unsure{section} $s$ de $\underline{S}^{n}(\underline{E})$ s'identifie à une \unsure{section} de $\underline{L}_{\underline{E}}$ sur $X=P(\underline{E})$, et \deleted{\unsure{a}} \inserted{définit} donc un \uline{diviseur} $\operatorname{div}(s)\deleted{\unsure{(s)}}$ sur $X$; la \unsure{condition} que $s$ [sur $Y$] ne \uline{s'annule pas}, signifie que ce diviseur induit bien un diviseur sur chaque fibre, \deleted{\unsure{D}} i.e. que $s$ ne s'annule pas.

D'ailleurs \deleted{\unsure{cette correspondance}} $\operatorname{div}(s)=\operatorname{div}(t)$ \deleted{\unsure{si et}} si et seulement si on \unsure{peut} trouver $\rho\in\Gamma(X,\underline{O}_X^*)$ t.q. $t=\rho s$, or $\rho$ s'identifie à un élément de $\Gamma(Y,\underline{O}_Y^*)$, (car $f_*(\unsure{\mathscr{O}|\underline{O}}_X)=\underline{O}_Y$). Donc $\operatorname{div}(s)=\operatorname{div}(t)\Leftrightarrow s$ et $t$ définissent la $=$ section de $\underline{S}^{n}(\underline{E})^*/\underline{O}_Y^*$.

De plus, $s\unsure{-|\to}\operatorname{div}(s)$ est compatible aux localisations sur $S$. D'où : \uline{on a une application injective} de \deleted{$\unsure{\Phi|\varphi}$} $\underline{\Phi}^{(n)}(X)(Y)=\Gamma(Y,\underline{S}^{n}(\underline{E})^*/\underline{O}_Y^*)$ \uline{dans l'ens des diviseurs loc. principaux \inserted{positifs} sur $X$ qui induisent \deleted{un diviseur} sur chaque fibre un diviseur, \inserted{qui est} de degré $n$}. Appelons de tels diviseurs sur $X$ des ``diviseurs transversaux de degré $n$''.

Je dis que l'application précédente est \uline{surjective}, i.e. que tout diviseur transversal \inserted{$D$} de degré $n$ provient d'une section de $\Phi^{(n)}(X)$. En effet, soit $\underline{L}=\underline{L}(D)$ le faisceau inversible défini par
\scanpage{60}
% archive sheet number: 59
$D$, \deleted{on sait que} la question étant locale \inserted{sur $Y$}, on peut supposer que $\underline{L}\simeq\underline{O}_X(k)$ par la \unsure{remarque}, d'autre part on \unsure{voit|sait} aussitôt, puisque $\underline{L}\mid f^{-1}(y)\simeq\underline{O}_{f^{-1}(y)}(k)$, que $k=n$. Donc $D$ est le diviseur d'une section de $\underline{O}_X(\unsure{n|k})$.

Ainsi, on a \deleted{\unsure{prouvé}} trouvé un isomorphisme \inserted{\unsure{ensembliste}}
\[
\underline{\unsure{\Phi|\varphi}}^{(n)}(P(\underline{E}))(Y)\xrightarrow{\sim}\unsure{\mathscr{D}|D}^{(n)}(X/Y).
\]
Cet isomorphisme est fonctoriel, compatible aux images inverses.

D'autre part, la notion de ``diviseur \inserted{loc. principal positif} transversal'' est susceptible de descente [car c'est la notion de : fibré inversible, muni d'une section régulière qui est non div. de $0$ et non diviseur de $0$ sur chaque fibre]. De plus, la notion de degré d'un diviseur sur un schéma algébrique de \unsure{Sévéri|Séveri}-Brauer a un sens \deleted{[\unsure{on étend la base} \inserted{\unsure{alg. clos}}]}, invariant par extension du corps de base, d'où la notion des \uline{diviseurs \inserted{positifs} transversaux de degré $n$} sur un \unsure{sch. Brauer}. Cette notion est susceptible de descente. D'où le

\uline{\unsure{Théorème}} Les sections de $\Phi^{(n)}(X)$ ($X$ un \unsure{Sév. Br.}) sont les diviseurs \inserted{positifs} transversaux de degré $n$ de $X/Y$.
\scanpage{61}
% archive sheet number: 60
\uline{Exemple de diviseur transversal}. Soit $X$ \deleted{de type fini} sur $Y$, \deleted{$\unsure{R}$} $R_{\unsure{0|o}}(X)$ le faisceau total des fractions de $\underline{O}_X$, \deleted{$\unsure{R_0}$} Alors $\Gamma(R_0(X)^*/\underline{O}_X^{\unsure{*|\bullet}})\simeq$ \deleted{\unsure{groupe}} \deleted{Ensemble des diviseurs} \inserted{idéaux fractionnaires} \unsure{divisoriels|divisibles} \inserted{loc.} principaux de $X\simeq$ diviseurs loc. principaux sur $X$. Soit $\underline{\unsure{J|I}}$ un idéal de $R_0(X)$, on donnera des conditions pour qu'il soit \deleted{\unsure{princ.}} \uline{divisoriel loc. principal transversal} \deleted{$\unsure{\xi}$}. On suppose $\underline{J}$ positif, i.e. $\underline{J}\subset\underline{O}_X$, et on suppose à priori \uline{même} que $\underline{J}$ est un idéal de $\underline{O}_X$, \deleted{$\unsure{\xi}$} donc correspond à un sous-schéma $X'$ de $X$.

\uline{Proposition} Supposons $X$ plat sur $Y$, $X,Y$ loc. noeth., soit $X'$ un sous-schéma fermé de $X$. On suppose $X$ \uline{plat} sur $Y$. Soit $x\in X'$; \inserted{$y$ son image sur $Y$.} Pour que $X'$ soit divisoriel \deleted{\unsure{dans un voisinage de $x$}} \inserted{et transversal $\unsure{\equiv|=}$ $Y$-transversal en $x$}, \unsure{il f.s.|il suffit} que l'on ait les conditions suivantes

\begin{enumerate}[label=(\roman*)]
\item $X'$ est $Y$-plat en $x$
\item $X'\otimes k(y)$ est divisoriel \deleted{\unsure{dans}} en $x$ dans $X\otimes k(y)$.
\end{enumerate}

\uline{Démonstration}. \uline{Nécessité} On peut supposer $\underline{O}_{X'}=\underline{O}_X/\unsure{\varphi|\phi}\underline{O}_X$, $\varphi\in\Gamma(X,\underline{O}_X)$, $\varphi$ \deleted{n'est pas diviseur de} \inserted{une fonction} \deleted{\unsure{de $0$}} \uline{\unsure{divisible|divisante}}, et $\varphi$ induisant sur les fibres des fonctions \unsure{divisibles|divisantes}. Alors on a une suite exacte
\[
0\unsure{-|\to}\underline{O}_X\xrightarrow{\varphi}\underline{O}_X\to\underline{O}_{X'}\to0,
\]
\deleted{\unsure{car $\varphi$}} \inserted{et si $z\in Y$} \deleted{\unsure{pour $\underline{O}_{X'}$ est plat}}
\[
0-\underline{O}_X\otimes k(z)\xrightarrow{\varphi_z}\underline{O}_X\otimes k(z)\to\underline{O}_{X'}\otimes k(z)\to\unsure{0}
\]
\scanpage{62}
% archive sheet number: 61
est exacte; donc, comme $\underline{O}_X$ est plat sur $Y$, il s'ensuit (réf) que $\underline{O}_{X'}$ est plat sur $Y$, de plus $\underline{O}_{X'}\otimes k(z)$ est \inserted{\unsure{un quotient|en général}} divisoriel pour $\underline{O}_X\otimes k(z)$. \inserted{Faisons donc $z=y$.}

\uline{Suffisance} La suite exacte $0\to\underline{\unsure{J|I}}\to\underline{O}_X\unsure{-|\to}\underline{O}_{X'}-0$ donne, puisque $\underline{O}_{X'}$ est $Y$-plat \uline{en $x$}, la suite exacte
\[
0\to\underline{J}\otimes k(y)\to\underline{O}_X\otimes k(y)\to\underline{O}_{X'}\otimes k(y)-0,
\]
et comme $\underline{O}_{X'}\otimes k(y)$ est divisoriel \inserted{en $x$}, il s'ensuit que $(\underline{J}\otimes k(y))_x$ est monogène, donc (Nakayama) $\underline{J}$ \deleted{\unsure{alors}} est \deleted{\unsure{engendré par un élém.}} \inserted{monogène} au voisinage de $x$, soit $\underline{J}=\unsure{\varphi|\phi}\underline{O}_X$ [en rapetissant $X$]. Comme $\varphi$ dans $\underline{O}_X\otimes k(y)$ est un générateur de $\underline{J}\otimes k(y)$, il n'est pas div. de $0$ dans $\underline{O}_X\otimes k(y)$, i.e. la suite
\[
0-\underline{O}_X\otimes k(y)\xrightarrow{\varphi}\underline{O}_X\otimes k(y)\to\underline{O}_{X'}\otimes k(y)\to0
\]
déduite de la suite
\[
0-\underline{O}_X\xrightarrow{\varphi}\underline{O}_X\to\underline{O}_{X'}\to0
\]
est \uline{exacte}. Je dis que $\varphi\colon\inserted{\underline{O}_X\to\underline{O}_X}$ est \uline{injectif} en $x$ : soit $I$ son noyau, on aura alors les suites exactes
\[
0-I\to\underline{O}_X\xrightarrow{\varphi}\deleted{\unsure{I}}\unsure{\underline{J}|J}\to0,\quad 0-\underline{J}\to\underline{O}_X\to\underline{O}_{X'}\to0,
\]
or $J$ est plat (car $\underline{O}_X$ et $\underline{O}_{X'}$ le sont) donc on a la suite exacte
\[
0-I\otimes k\to\underline{O}_X\otimes k\to J\otimes k\to0
\]
d'où $I\otimes k=0$ d'où $I=0$ (Nakayama). \unsure{c.q.f.d.}

\uline{Cor} Stabilité par extension de la base.
\scanpage{63}
% archive sheet number: 62
\uline{Groupe de Brauer global.}

Pour avoir des fondements solides et bien connus, nous nous plaçons en Géométrie \deleted{Algébrique} \inserted{analytique} [sur le corps $\mathbb{C}$ des complexes] Les groupes de cohomologie seront pris p.r. à la topologie usuelle.

Du \deleted{\unsure{les}} \inserted{diagramme de} suites exactes
\diagram{\begin{tikzcd}[ampersand replacement=\&,column sep=small]
e\arrow[r] \& \mathbb{C}^*\arrow[r] \& \operatorname{Gl}(n)\arrow[r] \& \operatorname{GP}(n-1)\arrow[r] \& e \\
e\arrow[r] \& \mathbb{Z}/n\mathbb{Z}\arrow[u]\arrow[r,dash] \& \operatorname{Sl}(n)\arrow[u]\arrow[r] \& \deleted{\unsure{\operatorname{Gl}}}\operatorname{GP}(n-1)\arrow[u]\arrow[r,dash] \& e
\end{tikzcd}}
on tire (pour un \unsure{bon} $X$) la ``suite exacte''
\diagram{\begin{tikzcd}[ampersand replacement=\&,column sep=tiny]
\cdots\arrow[r] \& H^1(X,\underline{\mathbb{C}}^*)\arrow[r] \& H^1(X,\underline{\operatorname{Gl}}(n))\arrow[r] \& H^1(X,\operatorname{GP}(n-1))\arrow[r] \& H^2(X,\mathbb{C}^*) \\
\cdots\arrow[r] \& H^1(X,\mathbb{Z}/n\mathbb{Z})\arrow[u]\arrow[r] \& H^1(X,\underline{\operatorname{Sl}}(n))\arrow[u]\arrow[r] \& H^1(X,\operatorname{GP}(n-1))\arrow[u]\arrow[r,dash] \& H^2(X,\mathbb{Z}/n\mathbb{Z})\arrow[u]
\end{tikzcd}}
Considérons, pour un \unsure{fibré} donné $\unsure{\omega|\sigma|\varpi}\in H^1(X,\operatorname{GP}(n-1))$, l'élément
\[
\delta(\omega)\in H^2(X,\mathbb{Z}/n\mathbb{Z})
\]
obstruction à remonter le groupe structural à $\operatorname{Sl}(n)$. Alors l'obstruction à remonter le groupe structural à $\operatorname{Gl}(n)$ est l'image de $\delta(\omega)$ dans $H^2(X,\mathbb{C}^*)$; elle est donc d'ordre $n$ [l'ensemble des élémts d'ordre $n$ dans $H^2(X,\mathbb{C}^*)$ étant exactement l'image de $H^2(X,\mathbb{Z}/n\mathbb{Z})$].

Les \unsure{classes} de fibrés sur $X$, de groupe $\operatorname{GP}(n-1)$, sont aussi les \unsure{classes} de fibrés en algèbres \deleted{\unsure{matrices}} de \deleted{\unsure{dim.}} \inserted{\unsure{r.g|r}} $n^2$; celles qui proviennent de $\operatorname{Gl}(n)$ sont les fibrés en algèbres qui
\scanpage{64}
% archive sheet number: 63
sont de la forme $L(E;E)$, où $E$ est un fibré vectoriel de dim. $n$.

Si $\unsure{A|\mathscr{A}},\unsure{B|\mathscr{B}}$ sont deux fibrés en algèbres de \deleted{rang} $m^2,n^2$, alors $A\otimes B$ est un fibré en algèbres de rang $m^2n^2=(mn)^2$, donc définit un élément de $H^1(X,\underline{\operatorname{GP}}(mn-1))$. \unsure{Ceci} est associé à une représentation de groupes
\[
\operatorname{GP}(m-1)\times\operatorname{GP}(n-1)\unsure{-|\to}\operatorname{GP}(mn-1)
\]
déduite de
\[
\operatorname{Gl}(m)\times\operatorname{Gl}(n)\longrightarrow\operatorname{Gl}(mn)
\]
par passage au quotient. Ainsi l'ens
\[
\unsure{\mathscr{B}|\mathcal{B}|B}(X)=\coprod_{n\geq\unsure{1|2}}H^1(X,\operatorname{GP}(n-1))
\]
est un monoïde \inserted{comm., \deleted{associatif}}, dont l'élémnt unité est l'unique élément de degré $1$. On \deleted{\unsure{met}} introduit dans ce monoïde une relation d'équivalence : deux \deleted{\unsure{fibrés}} élémts $a,b$ sont équivalents, si on peut trouver $c$, $c$ étant ``trivial'' --- i.e. provenant d'un fibré vectoriel --- t.q. $ac=bc$. Si $a\sim b$ et $a'\sim b'$, alors $aa'\sim bb'$, donc l'ensemble quotient, noté $\operatorname{Br}(X)$, est un monoïde \deleted{C'est un} (ass, comm. avec $1$). C'est \unsure{même} un groupe, car si
